\documentclass[11pt]{article}

\usepackage[
  letterpaper,
  left=0.70in,
  right=0.70in,
  top=0.75in,
  bottom=0.75in
]{geometry}
\usepackage[T1]{fontenc}
\usepackage[utf8]{inputenc}
\usepackage{mathptmx}
\usepackage{microtype}
\usepackage[table]{xcolor}
\usepackage{enumitem}
\usepackage{amsmath,amssymb}
\usepackage{booktabs}
\usepackage{tabularx}
\usepackage[colorlinks=true,allcolors=blue!55!black]{hyperref}

\definecolor{responseblue}{RGB}{0,70,140}
\definecolor{reviewerred}{RGB}{190,25,25}
\definecolor{placeholderred}{RGB}{155,25,25}
\definecolor{draftgray}{RGB}{246,246,246}

\newcommand{\manuscriptid}{26-3841v2}
\newcommand{\papertitle}{VLM-DexSafeDagger: VLM-Guided Intervention Policy Towards Safer Visuomotor Skill Distillation}
\newcommand{\placeholder}[1]{\textcolor{placeholderred}{\textbf{[#1]}}}
\newcommand{\response}{\par\smallskip\noindent\textcolor{responseblue}{\textbf{[Response:]}}\ }

\newcounter{reviewer}
\newcounter{reviewcomment}
\renewcommand{\theHreviewcomment}{\arabic{reviewer}.\arabic{reviewcomment}}

\newcommand{\reviewerheading}[1]{%
  \clearpage
  \stepcounter{reviewer}
  \setcounter{reviewcomment}{0}
  \phantomsection
  \pdfbookmark[1]{Response to #1}{response.reviewer.\arabic{reviewer}}
  \begin{center}
    {\Large\bfseries\scshape\color{reviewerred} Response to #1}\par
  \end{center}
  \vspace{0.25\baselineskip}
}

\newcommand{\reviewcommentlabel}[2]{%
  \refstepcounter{reviewcomment}
  \par\medskip
  \noindent\colorbox{draftgray}{%
    \parbox{\dimexpr\linewidth-2\fboxsep\relax}{%
      \textbf{#1.}\ \textbf{\itshape #2}%
    }%
  }\par\smallskip
}

\newenvironment{draftresponse}{%
  \response
  \textcolor{responseblue}{\textbf{Drafting checklist:}}
  \begin{itemize}[leftmargin=1.6em,itemsep=2pt,topsep=3pt]
}{%
  \end{itemize}
}

\newcommand{\draftitem}[1]{\item \placeholder{#1}}

\setlist[enumerate]{leftmargin=2em,itemsep=2pt,topsep=2pt}
\setlist[itemize]{leftmargin=1.6em,itemsep=2pt,topsep=2pt}
\setlength{\parindent}{0pt}
\setlength{\parskip}{3pt}
\setlength{\emergencystretch}{3em}
\linespread{1.12}

\begin{document}
\pagenumbering{roman}

\begin{center}
  {\Large\bfseries Authors' Responses}\par
  \vspace{0.35em}
  {\large\bfseries\color{placeholderred} WORKING SKELETON: NOT FOR SUBMISSION}\par
\end{center}

\vspace{1.5em}
\begin{tabularx}{\linewidth}{@{}p{0.28\linewidth}X@{}}
\textbf{Journal Title:} & IEEE Robotics and Automation Letters \\
\textbf{Reference Number:} & \manuscriptid \\
\textbf{Title of Paper:} & \papertitle \\
\textbf{Authors:} & Anonymous Authors
\end{tabularx}

\vspace{1.5em}
We sincerely thank the Associate Editor and the two anonymous reviewers for
their careful assessment of the revised manuscript. We appreciate their
recognition of the improvements made in the previous revision and their clear
identification of the remaining issues.

\placeholder{Replace this paragraph with a concise summary of the completed
revision after every response below has been factually verified. The summary
should cover teacher-reward reproducibility, baseline fairness, result
consistency, statistical analysis, the scope of the VLM contribution, and the
strengthened limitations statement.}

The detailed point-by-point responses follow. Reviewer comments are reproduced
in bold italic text. All red bracketed text is an internal drafting prompt and
must be replaced before submission.

\reviewerheading{the Associate Editor}

We thank the Associate Editor for summarizing the remaining concerns and for
directing us to the detailed reviews.

\reviewcommentlabel{Associate Editor Comment}{The revised manuscript has
improved considerably, and the reviewers acknowledge the authors' efforts in
addressing the previous comments. However, several important concerns remain,
particularly regarding the RL teacher-training procedure, the fairness of the
baseline comparison, inconsistencies in the reported results, and the evidence
supporting the contribution of VLM-guided adaptation. These issues need to be
addressed to properly support the claimed contributions. Please see the
individual reviews for further details.}

\begin{draftresponse}
  \draftitem{Summarize the exact teacher-reward equations, weights, scales, and
  truncation/bootstrap treatment added to the paper or supplement.}
  \draftitem{State the verified resolution of the DexTrAH/FABRICS comparison:
  either report a fair re-tuned baseline or withdraw the unsupported
  superiority claim and relabel the transferred implementation accurately.}
  \draftitem{Summarize the corrected Figure 3/Table III relationship and the
  unified definition of the reported means.}
  \draftitem{Summarize the final evidence and claim scope for VLM-guided
  adaptation without exceeding what the audited experiments establish.}
\end{draftresponse}

\reviewerheading{Reviewer \#29}

We sincerely thank the reviewer for recognizing the substantial improvements
in the revised manuscript and for identifying the remaining interpretive,
statistical, and reproducibility issues. Below are our point-by-point
responses.

\reviewcommentlabel{Comment 1}{The revision appropriately removes formal
``safety-critical'' claims. Please maintain this discipline throughout the
manuscript and consistently describe the method as empirical risk reduction
under task-specific unsafe-event definitions.}

\begin{draftresponse}
  \draftitem{Report the completed manuscript-wide audit of safety language,
  including the title, abstract, contributions, methods, captions, results,
  discussion, conclusion, supplement, and response.}
  \draftitem{State explicitly that the framework provides empirical reduction
  in task-specific unsafe-event rates and no formal safety guarantee.}
  \draftitem{List the principal locations where claims were revised.}
\end{draftresponse}

\reviewcommentlabel{Comment 2}{Real-world unsafe-event rates on unseen objects
remain around 39\%, which is substantial. This limitation should be emphasized
prominently in the abstract, discussion, and conclusion to avoid implying
deployment-ready safety.}

\begin{draftresponse}
  \draftitem{Verify and report the exact unseen-object unsafe-event rate,
  numerator, denominator, object count, and seed/trial aggregation.}
  \draftitem{Quote or summarize the new abstract, discussion, and conclusion
  language that foregrounds this limitation.}
  \draftitem{State plainly that the current system is not deployment-ready and
  identify the remaining sim-to-real/generalization gap.}
\end{draftresponse}

\reviewcommentlabel{Comment 3}{The real-robot difference between VLM-guided
and fixed-threshold critic policies appears modest. Report confidence intervals
and paired or proportion tests to determine whether the observed unsafe-rate
reductions are statistically distinguishable.}

\begin{draftresponse}
  \draftitem{Determine and state whether trials are genuinely paired by object,
  seed, initial pose, and other initial conditions.}
  \draftitem{Report the effect size for VLM-guided minus fixed-threshold
  unsafe-event rate, with confidence intervals for seen and unseen objects.}
  \draftitem{Name the justified statistical test, report its statistic and
  p-value, and state plainly whether the difference is statistically
  distinguishable.}
  \draftitem{Identify the table or paragraph containing the new analysis.}
\end{draftresponse}

\reviewcommentlabel{Comment 4}{The Bellman critic achieves a 5.8\%
false-negative rate on unsafe episodes. Please discuss the consequences of
these missed warnings and whether threshold adjustment can trade task success
for lower false-negative risk.}

\begin{draftresponse}
  \draftitem{Verify the 5.8\% numerator, denominator, and episode-level
  definition.}
  \draftitem{Explain that a missed critic warning can prevent critic-triggered
  takeover, while the disagreement branch may still trigger intervention.}
  \draftitem{If saved critic scores are available, report a threshold sweep
  covering false negatives, intervention rate, task success, and unsafe-event
  rate; otherwise disclose why the tradeoff cannot be reconstructed.}
\end{draftresponse}

\reviewcommentlabel{Comment 5}{Clarify whether critic metrics are computed at
the episode or time-step level and how early an unsafe event is detected before
termination. Warning lead time is important for intervention usefulness.}

\begin{draftresponse}
  \draftitem{State whether each critic metric is episode-level or time-step
  level and give the exact counting rule.}
  \draftitem{Define the first warning time, unsafe-event time, and lead-time
  computation at the 60 Hz control rate.}
  \draftitem{Report the verified lead-time distribution and useful-warning
  fractions, or state explicitly that saved time-step scores are insufficient
  to reconstruct them.}
\end{draftresponse}

\reviewcommentlabel{Comment 6}{The teacher policy is implicitly treated as a
safe fallback, yet teacher interventions can still generate unsafe contacts.
Report teacher unsafe-event rates and explain what happens when teacher and
critic recommendations conflict.}

\begin{draftresponse}
  \draftitem{Report teacher unsafe-event rates using an explicit, consistent
  aggregation population and denominator.}
  \draftitem{Clarify that teacher takeover is empirically safer in the tested
  setting but is not guaranteed safe.}
  \draftitem{Explain the conservative OR rule accurately: disagreement and
  critic output determine whether the gate requests teacher takeover; they are
  not two action recommendations.}
  \draftitem{State what occurs when the two gate signals disagree and note that
  teacher failures remain possible after takeover.}
\end{draftresponse}

\reviewcommentlabel{Comment 7}{Provide the final exact VLM identifier, prompt,
decoding settings, visual-sample selection rules, update interval, latency
distribution, failure handling, and cost in the main paper or permanent
supplementary material.}

\begin{draftresponse}
  \draftitem{Insert the final verified model identifier and provider/API
  version rather than a provisional model name.}
  \draftitem{Point to the exact prompt and machine-readable output schema.}
  \draftitem{Report decoding parameters, image resolution/detail settings,
  visual-sample selection rules, query schedule, parsing and clipping logic,
  and failure fallback.}
  \draftitem{Report latency distribution, token usage, and monetary cost from
  archived logs, distinguishing measured values from estimates.}
\end{draftresponse}

\reviewcommentlabel{Comment 8}{The ablations convincingly motivate
visual-semantic adaptation on representative training objects. Please state
more explicitly whether the VLM adaptation itself was evaluated on novel scene
semantics, rather than only novel object geometries.}

\begin{draftresponse}
  \draftitem{Audit the evaluation logs and define exactly which visual factors
  varied: object geometry, texture, material, lighting, background, camera, or
  semantic hazard type.}
  \draftitem{State explicitly whether novel scene semantics were evaluated.}
  \draftitem{If not, narrow the contribution claim and identify semantic
  generalization as future work rather than inferring it from novel geometry.}
\end{draftresponse}

\reviewcommentlabel{Comment 9}{The current unsafe-event taxonomy is
task-specific. Clarify which events are directly sensor-observable on hardware
and which depend on privileged simulation signals during training or
evaluation.}

\begin{draftresponse}
  \draftitem{Create an event-by-event mapping for harmful collision, object
  escape, palm flip, and hand-too-far signals.}
  \draftitem{For each event, distinguish hardware sensor or vision observables
  from simulator-only state, contact, force, or geometry information.}
  \draftitem{Explain how real-robot events were labeled and whether any
  post-processing or manual annotation was used.}
\end{draftresponse}

\reviewcommentlabel{Comment 10}{Please add a compact limitations paragraph
covering the remaining sim-to-real contact-strategy mismatch, external VLM
dependence, task-specific safety definitions, and absence of formal constraint
guarantees. No major methodological restructuring appears necessary.}

\begin{draftresponse}
  \draftitem{Insert the final limitations paragraph or quote it here.}
  \draftitem{Confirm that it covers all four requested limitations without
  weakening or omitting the 39\% unseen-object unsafe-event limitation.}
  \draftitem{Identify its location in the manuscript.}
\end{draftresponse}

\reviewerheading{Reviewer \#31}

We sincerely thank the reviewer for the detailed technical assessment, for
recognizing the expanded real-robot evaluation and reproducibility material,
and for identifying the remaining issues in teacher training, baseline
fairness, result consistency, and the scope of the VLM contribution.

\reviewcommentlabel{Major Comment M1}{The reward for teacher-policy training is
described only verbally. Table I lists only the purpose of each reward term;
the exact formulation, weights, and scales of the terms are not given in the
manuscript. Since teacher training is the starting point of every subsequent
stage, the work does not appear reproducible in its current form. Please
provide the exact formulation and weights of each reward term.}

\response We thank the reviewer for identifying this reproducibility gap.  We
agree that the verbal descriptions in Table~I were insufficient to reproduce
the teacher-training objective.  We have therefore added a new
``Teacher-Reward Formulation'' subsection to Supplementary Section~A.2.  The
new subsection first states the complete additive teacher reward and then gives
the exact mathematical expression for each of its eight terms:
$r_{\mathrm{dist}}$, $r_{\mathrm{curl}}$, $r_{\mathrm{align}}$,
$r_{\mathrm{contact}}$, $r_{\mathrm{near}}$, $r_{\mathrm{goal}}$,
$r_{\mathrm{lift}}$, and $r_{\mathrm{action}}$.  It also defines the geometric
distances, contact and grasp indicators, action decomposition, activation
conditions, and clipping operations used by these terms.

Supplementary Table~2 collects all numerical coefficients and implementation
parameters, including the distance and contact thresholds, reference finger
posture, ADR schedule, lift delay, relative hand-action penalty, and global PPO
reward scaling.  We verified every equation and value directly against the
implementation and configuration used to train the reported teacher policies.
This addition provides the requested reproducibility details while retaining
the compact qualitative overview in the main-paper Table~I.

\reviewcommentlabel{Major Comment M2}{The claimed superiority of the authors'
RL recipe is not supported. (1) The authors explain the performance gap between
ablations (a) and (b) by stating that the FABRICS controller was tuned for the
original platform and transfers poorly to the new hardware. This implies that
the DexTrAH baseline was not tuned for the authors' hardware, and therefore the
comparison cannot serve as evidence that the authors' recipe produces more
effective teacher policies. A comparison against a DexTrAH baseline re-tuned
for the authors' platform, or additional experiments and analysis that can
demonstrate the superiority of the recipe, is required. (2) The difference
between the ``DexTrAH-RGB-style RL recipe'' and ``our RL recipe'' is not defined
in the manuscript. Please state explicitly how the two recipes differ. (3) The
second contribution claims ``a safety-aware RL reward design that explicitly
accounts for risky events.'' However, the reward terms in Table I are generic
shaped-reward terms, and no term explicitly penalizes the unsafe events defined
by the authors, such as collisions or object escape. Early termination does
penalize unsafe events implicitly by truncating the future return, but it is
indistinguishable from a timeout unless truncation is bootstrapped. The
manuscript does not specify any of these. Please either add an ablation of the
reward terms or of the termination conditions, or moderate the contribution
statement.}

\response
\textcolor{responseblue}{\textbf{Drafting checklist by subpart:}}
\begin{enumerate}[label=(\arabic*),leftmargin=2em,itemsep=4pt,topsep=3pt]
  \item \placeholder{Choose and document the factual resolution: report a
  re-tuned embodiment-matched DexTrAH/FABRICS baseline, or withdraw the
  superiority claim and state that the transferred baseline is not evidence of
  recipe superiority.}
  \item \placeholder{Provide a precise component-by-component comparison of
  the two recipes, including embodiment, observations, action space, reward,
  termination, optimizer, network, training budget, and controller tuning.}
  \item \placeholder{Report the verified termination/bootstrap semantics and
  either add the requested ablation or moderate the reward-design contribution
  so it matches the actual implementation and evidence.}
\end{enumerate}

\reviewcommentlabel{Major Comment M3}{The unsafe rates in Fig. 3 and Table III
are inconsistent. In the ``Comparison Unsafe Episode Rate'' panel of Fig. 3,
(c) is the highest, (b) and (e) are nearly indistinguishable, and (f) is the
lowest. In contrast, Table III reports (b) 17.4 $>$ (c) 16.7 $>$ (e) 11.9 $>$
(f) 9.3, which differs in both values and ordering. Please revise so that the
two results are consistent with each other.}

\begin{draftresponse}
  \draftitem{Verify both outputs from raw logs and record objects, seeds,
  rollout phase, policy checkpoint, intervention state, episode denominator,
  smoothing/window, and aggregation population.}
  \draftitem{If Figure 3 is online-distillation data and Table III is frozen
  post-training evaluation, state this distinction prominently and explain why
  values/orderings need not match.}
  \draftitem{If they were intended to show the same measurement, identify and
  correct the plotting, transcription, or aggregation error.}
  \draftitem{Point to the revised captions and text.}
\end{draftresponse}

\reviewcommentlabel{Major Comment M4}{The contribution of the VLM is not
established by the (e) vs. (f) comparison alone. In particular, the third
contribution claims effectiveness ``across different visual contexts and unsafe
scenarios,'' but the results currently presented do not clearly verify this
claim. Please either moderate the corresponding contribution statement or
provide additional analysis.}

\begin{draftresponse}
  \draftitem{Choose the evidence-based route after auditing stored data:
  provide stratified visual-context/unsafe-scenario analysis, or narrow the
  contribution statement.}
  \draftitem{Do not treat object geometry alone as evidence of novel visual
  semantics.}
  \draftitem{Report the final claim language and the exact analysis, table, or
  limitation that supports it.}
\end{draftresponse}

\reviewcommentlabel{Minor Comment m1}{One paragraph in Sec. III-A has a line
spacing that differs from the rest of the text. Please check the formatting of
the entire manuscript.}

\begin{draftresponse}
  \draftitem{Identify and fix the source of the spacing inconsistency.}
  \draftitem{Report a full rendered-page formatting audit, including captions,
  tables, equations, section transitions, and line spacing.}
\end{draftresponse}

\reviewcommentlabel{Minor Comment m2}{The ``Mean'' columns of Table II and
Table III appear to average over different populations. Please unify the
definition of ``Mean'' across the two tables. In addition, please briefly state
why these four objects were selected from the 40-object set.}

\begin{draftresponse}
  \draftitem{Verify whether each current mean is macro-averaged or pooled and
  whether it covers four displayed objects or all 40 objects.}
  \draftitem{Use one explicit population definition in both tables, or label
  distinct populations unambiguously if unification is impossible.}
  \draftitem{Report the true four-object selection procedure, including whether
  it was fixed before inspecting comparative results. Do not invent a
  representativeness rationale after the fact.}
  \draftitem{Point to updated captions and any full 40-object results.}
\end{draftresponse}

\reviewcommentlabel{Minor Comment m3}{Fig. 4 is difficult to read at print
quality. Please also explain why only two objects (windcart and cow) are shown.}

\begin{draftresponse}
  \draftitem{Describe the figure redesign and verify legibility at final print
  size.}
  \draftitem{State the true selection rationale for WindCart and Cow and whether
  it was fixed before inspecting results.}
  \draftitem{If space permits, point to higher-resolution or complete
  per-object failure-mode results in permanent supplementary material.}
\end{draftresponse}

\vspace{1em}
We thank the Associate Editor and the reviewers again for their constructive
feedback. \placeholder{Replace this closing with a concise statement of the
verified revisions and their effect on the manuscript's claims.}

\end{document}
